\chapter{Reglaje, tracking y balanceo}
\label{cap:mm-reglaje}

\opening{Ángulos de pala medidos al 75\,\% del radio, tubos ajustados por medias vueltas y vibraciones llevadas por debajo de 0,2 pulgadas por segundo.}

\section{Introducción}
\fuente{MM, capítulo 10}

Este capítulo reúne el reglaje de los controles del rotor principal, del rotor de cola y la correlación del acelerador. Los procedimientos de \emph{tracking} y balanceo se usan junto con las instrucciones del equipo de balanceo (Chadwick-Helmuth Vibrex, TEC ACES, Dynamic Solutions MicroVib o equivalente). El ajuste de las rótulas de los tubos se hace según \S\,23-34 (capítulo~\ref{cap:mm-practicas}): tras cualquier ajuste, verifique por el agujero testigo el enganche de la rosca y apriete y marque contratuercas y palnuts.

\section{Controles del rotor principal}

\subsection{Cíclico}

El recorrido del cíclico no es ajustable: lo limita la placa de topes A211-1 en la caja del cíclico. Si cambió la longitud del tubo A121-1 o de la horquilla A205, reajústelos a las medidas del capítulo~\ref{cap:mm-controles}.

\begin{pasos}
\item Cíclico contra el tope trasero y colectivo abajo: ajuste el tubo A121-3 para \textbf{0,130~in} entre el brazo trasero del eje intermedio y la brida de la tapa superior de la caja principal.
\item Colectivo arriba y cíclico adelante a fondo: holgura mínima de \textbf{0,125~in} entre el brazo delantero del eje intermedio y la guía del tubo A121-7.
\item Lleve el cíclico a \textbf{neutro}: 8,3~in a la derecha del tope izquierdo y en el punto medio del recorrido longitudinal; colectivo abajo; aplique fricción total en ambos sin mover el cíclico.
\end{pasos}

\subsection{Plato oscilante}

Con los mandos bloqueados así, ajuste los tubos A121-7 y A121-5 para una holgura \textbf{constante} del plato inferior a la brida A281-1: mínimo \textbf{5/8~in}, nominal \textbf{3/4~in} (bloques de reglaje MT146-1).

\subsection{Colectivo}

El tope del deslizador no es ajustable; se verifica que haya recorrido completo sin interferir con el plato. Levante la funda y suba el colectivo a fondo: la parte superior de la bola (\emph{uniball}) debe quedar \textbf{al ras hasta 0,060~in por encima} del tubo deslizante. Si sobresale más, enrosque por igual las rótulas superiores de los A121-5 y A121-7 para bajar el plato.

\subsection{Preparación para medir ángulos de pala}

Los ángulos se miden al 0,75 del radio, \textbf{37,75~in hacia adentro desde la punta}, con el accesorio MT050-1 y un transportador de hélices Kell-Strom KS113 o equivalente.

\begin{pasos}
\item Todas las rótulas deben cumplir las tolerancias de juego radial y axial.
\item Reglaje del plato oscilante hecho.
\item Bieletas de paso inicialmente a \textbf{6,82~in} entre centros (la longitud final la da el ajuste de RPM de autorrotación).
\item Horquillas A205 del plato: inferior \textbf{3,50 $\pm$ 0,03~in}, superior \textbf{3,70 $\pm$ 0,03~in} entre centros.
\item Nivele el helicóptero sobre el cubo del rotor (\S\,18-12).
\item Con una vara de \emph{tracking}, bascule el rotor hasta un track de \textbf{$\pm$1~in}.
\item Marque con cinta la punta de una pala. Ponga el transportador sobre el cubo, paralelo al perno de balanceo y mirando a la pala marcada, y póngalo a cero. \textbf{Al medir, el transportador siempre debe mirar a la punta marcada} (si no, se duplica el error del instrumento).
\item Marque las palas a 37,75~in de la punta, una roja y otra azul.
\end{pasos}

\begin{advertencia}
Marque las palas solo con cinta, lápiz graso o marcador blando. Un bolígrafo u objeto punzante puede rayar el revestimiento y provocar fisuras y falla por fatiga.
\end{advertencia}

\subsection{Recorrido del cíclico y del colectivo}

En cada posición se mide el ángulo de cada pala con la bieleta en un lado, se gira el rotor 180° y se mide de nuevo; se promedian ambas lecturas.

\begin{center}
\small
\begin{tabularx}{\linewidth}{@{}>{\raggedright\arraybackslash}p{0.30\linewidth}l>{\raggedright\arraybackslash}X@{}}
\toprule
Posición & Ángulo promedio & Ajuste \\
\midrule
Cíclico adelante (palancas A958-1) & 10,5°--11,0° & A121-5 trasero (1 vuelta = 0,42°) o ambos A121-7 por igual \\
Cíclico adelante (palancas A175-1) & 8,3°--8,8° & ídem \\
Cíclico atrás & 8,5°--9,5° & ídem; si se ajusta, revisar adelante \\
Cíclico izquierda & 9,0°--9,5° & A121-7 derecho o izquierdo (1 vuelta = 0,6°) \\
Cíclico derecha & 5,5°--6,0° & ídem; si se ajusta, revisar izquierda \\
Colectivo abajo & 2° $\pm$ 0,5° cabeza arriba & bieletas; palas dentro de 0,2° entre sí \\
Colectivo arriba & 14° $\pm$ 0,5° cabeza arriba & --- \\
\bottomrule
\end{tabularx}
\end{center}

Para el cíclico longitudinal, las bieletas se alinean con el eje longitudinal; para el lateral, con el transversal; el cíclico se sujeta contra el tope (pueden usarse bolsas de arena). El ángulo final con colectivo abajo lo determina la RPM de autorrotación. Al final se verifica otra vez la posición de la bola del plato y se hace el tracking y balanceo.

\section{Controles del rotor de cola}

\begin{itemize}
\item \textbf{Pedales}: pasador de reglaje de 3/16~in (sirve un perno AN3 o NAS6603 largo) por el panel de quilla derecho y la palanca A317-1; ajuste los tubos A121-9 y -11 para \textbf{2,90 $\pm$ 0,03~in} de cada pedal a su tope en la consola.
\item \textbf{Palanca A316-1}: pedal izquierdo contra el tope delantero; ajuste el tubo vertical A121-13 para \textbf{0,060~in} mínimo entre el brazo de la palanca y el túnel del cortafuego vertical.
\item \textbf{Palanca A331-1}: pedal derecho contra el tope; ajuste el tubo horizontal A121-15 para \textbf{0,100~in} mínimo entre la palanca y la carcasa de la caja del actuador (con el actuador desacoplado). Con pedal izquierdo, verifique holgura con el cono de cola.
\item \textbf{Palanca A120-1}: pedales en neutro; ajuste el tubo A121-17 para una cota nominal de \textbf{4,85~in} entre el eje del brazo y la cara mecanizada de la pieza fundida del cono de cola.
\item \textbf{Bieletas}: 2,360~in entre centros, si son ajustables.
\end{itemize}

\begin{advertencia}
Las dos bieletas del rotor de cola deben tener el mismo número de parte; no mezcle de una pieza con ajustables, ni de acero con de aluminio.
\end{advertencia}

\subsection{Ángulos del rotor de cola}

\begin{pasos}
\item Con los pedales desmontables instalados, nivele sobre el cubo del rotor principal. Gire el rotor de cola hasta que la pala delantera quede paralela al cono y pegue una vara de tracking (chapa de 1$\times$12~in doblada a 90° a 2~in del extremo) junto a su punta.
\item Pedal izquierdo contra el tope: marque el paso del agujero de drenaje de cada punta (girando por el eje del embrague para no alterar el basculamiento) y bascule el rotor a la posición media; repita hasta un track dentro de \textbf{0,125~in}.
\item Marque cada pala a \textbf{5,25~in} de la punta (0,75 R). Con la pala delantera sujeta en la marca, ponga el accesorio MT050-2 centrado en la marca de la pala trasera (lado interior) y mida; gire 180° y mida la otra.
\item \textbf{Pedal izquierdo a fondo: 19,0°--19,5°} (punta a la derecha), ajustando las rótulas del tubo A121-17 (1 vuelta = 0,42°).
\item \textbf{Pedal derecho a fondo: 9,6°--10,6°} (punta a la izquierda), con nueva marca de tracking. Si se ajusta, revise el izquierdo.
\item Si no se logran ambos rangos, la suma de los dos ángulos indica el recorrido de pedales: \textbf{menos de 28,6°}, recorrido insuficiente (alargue los A121-9 y -11); \textbf{más de 30,1°}, excesivo (acórtelos), en medias vueltas.
\end{pasos}

\begin{precaucion}
La diferencia entre los ángulos de las palas azul y roja no debe superar \textbf{0,4°}; si no se logra con las bieletas, reemplácelas. Las contratuercas de las bieletas deben estar apretadas para que la medición sea exacta.
\end{precaucion}

\section{Correlación del acelerador}

\widefigure[0.85\linewidth]{assets/mm/acelerador-correlacion.png}{Reglaje de correlación del acelerador: resorte de sobrecarrera A327-1, brazo A609-2 (2,20 $\pm$ 0,03~in con O-360; 2,05 $\pm$ 0,03~in con O-320), tubo A336 o B364 y palanca del carburador (58° $\pm$ 7,5° a plena apertura). Fuente: MM, figura 10-8.\label{fig:acel-correlacion}}

\textbf{Verificación en servicio}: gire el acelerador hacia ``cerrado'' venciendo el resorte de sobrecarrera hasta el tope y, sosteniéndolo, suba el colectivo al tope superior: la palanca del carburador debe \emph{apenas empezar a moverse} al llegar arriba.

\textbf{Ajuste}: antes de tocar la biela del acelerador, el ralentí y el suplementado del motor deben estar correctos. Ajuste el resorte A327-1 sin compresión para la geometría de la figura (colectivo arriba, acelerador cerrado, brazo A609-2 a la cota indicada); luego ajuste el tubo A336 o B364 para que la palanca del carburador quede \textbf{0,020--0,050~in} separada del tornillo de ralentí (el tubo se reajusta siempre que se ajuste el tope de ralentí; pueden usarse hasta 2 separadores A130-54 con el B364). Conexión al agujero medio. Apriete contratuercas, verifique testigos y frene con alambre; con rótula A933-3 se requiere palnut B330-7.

\section{Tracking y balanceo}

\begin{nota}
Calibre el equipo según su fabricante, \textbf{al menos una vez al año} o si se cae, se usa mal o hay dudas.
\end{nota}

\subsection{Instalación del equipo}

\begin{itemize}
\item \textbf{Rotor principal}: acelerómetro con soportes en el panel A359-1 (se retiran los mandos desmontables); soporte del captador magnético PT121 en el plato oscilante, pernos a 100~in-lb + arrastre (queda instalado para el balanceo en vuelo); captador 3030 con separación del interruptor de \textbf{0,030 $\pm$ 0,010~in}; cables asegurados con cinta cada 12~in, sin posibilidad de enredarse con los pedales; cintas reflectantes 4270 en las palas. El interruptor debe estar \textbf{opuesto al brazo de cuerda} (si se instala del mismo lado, la carta queda desfasada).
\item \textbf{Rotor de cola}: acelerómetro 4177 bajo el perno superior delantero del cartucho de salida de la caja de cola, conector hacia arriba; cable enrollado al cono; cinta reflectante en la cara interior del cubo a 1~in del perno exterior de una pala; fotocelda alineada con la cinta a las 12.
\end{itemize}

\begin{precaucion}
La seguridad de los cables es esencial: el helicóptero se volará hasta $V_{NE}$. \textbf{El equipo de balanceo del rotor de cola debe retirarse antes del vuelo.}
\end{precaucion}

\subsection{Procedimiento del rotor principal}

Antes de instalar el equipo, verifique palas limpias y lisas; juegos de rótulas, rodamientos esféricos y tijeras dentro de límites; fricción correcta de la tijera superior, de inclinación del plato y de las bisagras de conicidad.

\begin{pasos}
\item Track en vuelo estacionario: acorte la bieleta de la pala alta hasta \textbf{0,25~in}.
\item Balanceo en estacionario según la carta, hasta \textbf{0,2~IPS} (pulgadas por segundo).
\item Track en vuelo a 50, 60, 70, 80, 90 y 100~nudos (sin exceder $V_{NE}$).
\item Corrija una pala que ``sube'' con la velocidad doblando su aleta compensadora hacia abajo.
\item Repita hasta un track dentro de \textbf{3/8~in} a todas las velocidades.
\item Reajuste el balanceo en estacionario a 0,2~IPS o menos.
\item Verifique y ajuste la RPM de autorrotación.
\item Evalúe la compensación del colectivo y del cíclico longitudinal y lateral.
\item Verifique de nuevo el balanceo en estacionario ($\le$0,2~IPS).
\end{pasos}

\begin{advertencia}
Un rotor que queda suave tras el balanceo pero se desbalancea a los pocos vuelos es sospechoso y debe ser examinado por RHC antes de volver a volar.
\end{advertencia}

\subsection{Correcciones de balanceo del rotor principal}

\begin{itemize}
\item \textbf{En envergadura}: pesos de punta bajo la tapa (las arandelas pueden recortarse). Tornillos de pesos y de tapa con antiagarrotante A257-9 y torque especial de \textbf{40~in-lb}. Equivalencias: 1 arandela AN960-10 (o NAS1149F0363P) = 2 AN960-10L (o NAS1149F0332P); A298-2 = 0,034~lb; A298-3 = 0,009~lb.
\item \textbf{En cuerda, grueso}: pasar suplementos del perno de balanceo de un lado al otro (bajo la cabeza o bajo la tuerca), desplazando el cubo.
\item \textbf{En cuerda, fino}: pesos del brazo de cuerda; total máximo \textbf{dos A255-2} o equivalente. 1 A255-1 = 8 arandelas AN970-4; 1 A255-2 = 3 A255-1.
\end{itemize}

\subsection{Ajuste de las bieletas de paso}

Para el track en estacionario se \textbf{acorta la pala alta}; en ajustes colectivos ambas bieletas se ajustan exactamente igual.

\begin{center}
\small
\begin{tabularx}{\linewidth}{@{}>{\raggedright\arraybackslash}p{0.32\linewidth}>{\raggedright\arraybackslash}p{0.3\linewidth}>{\raggedright\arraybackslash}X@{}}
\toprule
Ajuste (media vuelta) & Track & Ángulo \\
\midrule
Fino: barril (A258-5) o racor (A258-1) & 0,125~in & 0,12° \\
Medio: rótula superior & 0,25~in & 0,24° \\
Grueso: rótula inferior (A258-5) o conjunto de enlace (A258-1) & --- & 0,36° \\
\bottomrule
\end{tabularx}
\end{center}

Tras cualquier ajuste, la espiga roscada de la rótula debe \textbf{impedir el paso de un alambre de 0,020~in} por el agujero testigo. Tuerca 21FKF-616 con torque especial usando llave de contrafuerza; contratuerca y palnut superiores con torque estándar; en la A258-1, frenar racor y enlace con alambre de 0,032~in. Oriente las rótulas para permitir la máxima rotación de la bieleta sin trabarse.

\subsection{Aletas compensadoras}

Una pala que ``sube'' (separación mayor de 3/8~in con velocidad) se corrige \textbf{doblando hacia abajo la aleta de la pala alta}. Se marca una línea en la aleta en prolongación del borde de fuga, se pone a cero el comparador MT352-1 y se dobla con la herramienta MT526 apretando su perno (nunca usarla como palanca). \textbf{0,015~in hacia abajo en la aleta $\approx$ 0,2~in hacia abajo en la punta.}

\begin{precaucion}
\begin{itemize}
\item No use herramientas de otros fabricantes: dañan las palas Robinson.
\item MT526-1 es para palas A016-4 (revestimiento de acero) y MT526-8 para A016-6 (aluminio); no son intercambiables. MT090-1 es obsoleta.
\item El comparador MT352-1 debe tocar el borde de fuga de la aleta; si no lo toca, la aleta está doblada más allá del límite.
\item Doble hacia arriba solo si es imprescindible: puede aumentar la vibración.
\end{itemize}
\end{precaucion}

\subsection{Diagnóstico de tracking y balanceo}

\begin{center}
\small
\begin{tabularx}{\linewidth}{@{}>{\raggedright\arraybackslash}p{0.22\linewidth}>{\raggedright\arraybackslash}p{0.36\linewidth}>{\raggedright\arraybackslash}X@{}}
\toprule
Síntoma & Causa probable & Corrección \\
\midrule
Sacudida excesiva del cíclico & Rotor fuera de track o balanceo; horquilla A205-3 áspera o trabada; rodamientos de eje de pala con \emph{brinelling}; pintura saltada; rótulas de bieletas ásperas; funda de pala desalineada; bisagras trabadas; aletas dobladas hacia arriba; palas no emparejadas & Track y balanceo; reemplazo o reparación; repintar; doblar aletas hacia abajo por igual; RHC \\
Vibración excesiva de la aeronave & Track o balanceo; fricción de bisagras mal ajustada; rodamientos de conicidad pegajosos; \emph{brinelling}; rodamiento de balanceo gastado & Track y balanceo; ajuste de fricción; RHC \\
Fuerzas excesivas en el cíclico & \emph{Brinelling}; aletas dobladas hacia arriba & RHC; doblar aletas hacia abajo \\
Track intermitente & Fricción de bisagras; rodamientos pegajosos; bisagra de balanceo sin asentar; \emph{brinelling} & Ajustes; minimizar error de track; RHC \\
Cambios bruscos de compensación & Rodamientos de balanceo gastados; \emph{brinelling} & RHC \\
Vibración lateral intermitente & Fallo de encendido (bujías, cables, magneto) o motor sin asentar & Instrucciones de Lycoming \\
\bottomrule
\end{tabularx}
\end{center}

\subsection{Balanceo del rotor de cola}

Con el equipo instalado, seleccione en el balanceador el canal correspondiente, rango de RPM $\times$10 y dial en 340; con el helicóptero en marcha y \textbf{el gobernador conectado}, observe el rotor de cola con el Strobex, sintonice y registre ángulo de reloj e IPS. Ajuste hasta \textbf{menos de 0,2~IPS}.

\begin{itemize}
\item \textbf{Envergadura}: palas de punta cuadrada A029-1, pesos bajo la tapa de punta (A134-1/-2 o arandelas AN960-8/-8L; las -8L pueden recortarse). Palas de punta redonda A029-2, arandelas de distinto diámetro bajo la tuerca del perno exterior de la pala; deben sobresalir \textbf{2 a 4 hilos} tras el torque.
\item \textbf{Cuerda}: arandelas A141-14, A214-3, AN960-416 o -416L bajo la tuerca del perno de la bieleta, cambiando el largo del perno si hace falta para el enganche de rosca. Con A029-2 debe haber \textbf{al menos una arandela NAS1149F0432P} en la posición de cuerda de cada pala.
\item Equivalencias (A029-2): 2 NAS1149F0432P = 1 F0463P; 3,5 F0432P = 1 A214-3; 3 F0463P = 1 A141-14; 6,5 F0463P = 1 AN970-4.
\end{itemize}

\section{RPM de autorrotación}

\begin{advertencia}
Una RPM de autorrotación mal ajustada (demasiado baja) puede impedir alcanzar la RPM adecuada con pesos bajos. Durante la prueba no permita la sobrevelocidad del rotor: alargue progresivamente ambas bieletas hasta poder bajar el colectivo a fondo sin sobrevelocidad.
\end{advertencia}

\begin{pasos}
\item Haga la prueba con el menor peso práctico; registre peso de despegue y horómetro.
\item Altímetro en 29,92~inHg (1013,2~hPa). Autorrote con el colectivo firmemente contra el tope inferior a \textbf{50~KIAS}.
\item Tome al menos 3 lecturas a intervalos de 500 a 1000~ft, registrando horómetro, altitud de presión, OAT y \% RPM.
\item Para cada lectura, peso de prueba = peso de despegue $-$ combustible consumido (tiempo transcurrido $\times$ \textbf{45~lb/h}); obtenga la RPM de la gráfica (OAT $\rightarrow$ altitud de presión $\rightarrow$ peso $\rightarrow$ \% RPM) y la corrección = RPM de gráfica $-$ RPM de prueba.
\item Ajuste según la corrección promedio: \textbf{alargar ambas bieletas baja la RPM} (una vuelta completa de rótula $\approx$ 3\,\%; una vuelta de rótula = dos vueltas del barril); acortarlas la sube. Ajuste ambas exactamente igual para no alterar el track.
\item Repita hasta quedar dentro de \textbf{$\pm$1\,\%} de la RPM de la gráfica.
\end{pasos}

\widefigure[1.0\linewidth]{assets/mm/rpm-autorrotacion.png}{RPM de autorrotación con colectivo a fondo y 50~KIAS en función de temperatura, altitud de presión y peso bruto (ejemplo: OAT 23~°C, Hp 2000~ft, 1050~lb $\rightarrow$ 99,5\,\%). Fuente: MM, figura 10-15.\label{fig:rpm-autorrot}}
